\section{Ajakava}

\begin{table}[H]
    \centering
    \noindent\setlength\tabcolsep{4pt}%
    \begin{tabularx}{\linewidth}{|l|>{\RaggedRight\arraybackslash}X|}
        \hline
        \textbf{Daatum} & \textbf{Tegevus}\\
        \hline
        \hline
        september 2024 & Projekti eeltutvustus, meeskonna liidri CanSat NeXTi komplektiga tutvumine\\
        \hline
        oktoober 2024 & Liikmete leidmine ning meeskonna loomine\\
        \hline
        oktoober –-- november 2024 & CanSat projekti teistele tutvustamine ning rollide jagamine\\
        \hline
        november --- detsember 2024 & Õppematerjalidega tutvumine, vajalike andurite tellimine ning materjalide ostmine\\
        \hline
        jaanuar --- veebruar 2025 & Töö projektiga, prototüübi (algmudeli) valmistamine, esimesed katsed\\
        \hline
        märts 2025 & Valmismudeli testimine ning muutmine, aruanne koostamine\\
        \hline
        aprill 2025 & Viimased ettevalmistused võistlusteks\\
        \hline
        18. aprill 2025 & Kogu konstruktsiooni testimine, alla heitmine Sinimägede vaatetornilt\\
        \hline
        18.--23. aprill 2025 & Viimased parandamised ja seadistamised\\
        \hline
        24. aprill 2025 & VÕISTLUS\\
        \hline
    \end{tabularx}
    \caption{Ajakava}
    \label{tab:2}
\end{table}

Projekti algusest kuni talvekoolivaheajani kohtusime juhendajaga kord kuus, alates märtsist --- kord nädalas.

\section{Ressursid}

Projekti eelarve nõuded: 
Maksimaalne eelarve: Lõpliku CanSat-mudeli kogueelarve ei tohi ületada 500 eurot. Kui kasutatakse CanSat NeXT komplekti, arvestatakse selle väärtuseks 70 eurot.

\begin{table}[H]
    \centering
    \noindent\setlength\tabcolsep{4pt}%
    \begin{tabularx}{\linewidth}{|>{\RaggedRight\arraybackslash}X|l|}
        \hline
        \textbf{Nimetus} & \textbf{Hind (euro)}\\
        \hline
        \hline
        CanSat NeXT komplekt & 70\\
        \hline
        Õhuniiskuse ja CO$_2$ andur SCD41& 20\\
        \hline
        Juhtmed & 2\\
        \hline
        Pintsetid & 5\\
        \hline
        Langevarju materjal (silk, niidid, nailon, nöörid jne) & 40\\
        \hline
        Plastmassist kinnitused & 2,40\\
        \hline
        Karabiinid & 2\\
        \hline
        Korpuse 3D printimine & 20\\
        \hline
        Jootmise materjalid, kinnitused, kruvid jne & 10\\
        \hline
        Pehmendusmaterjal & 10\\
        \hline
        \textbf{KOKKU} & 181,40\\
        \hline
    \end{tabularx}
    \caption{Kasutatud materjalid}
    \label{tab:3}
\end{table}

\section{Riskianalüüs (riskid ning nende vähendamine + riskid stardipäeval)}

Riskid, mis olid ning võivad olla projekti realiseerimise käigus ning nende vähendamise meetmed esitatud tabelina:

\begin{table}[H]
    \centering
    \noindent\setlength\tabcolsep{4pt}%
    \begin{tabularx}{\linewidth}{|l|>{\RaggedRight\arraybackslash}X|}
        \hline
        \textbf{Risk} & \textbf{Riski vähendamise meetod}\\
        \hline
        \hline
        Mitte jõuda tähtajaks & Jälgida koostatud ajakava\\
        \hline
        Tehniline osa ei tööta & Katsetada saatja ja vastuvõtja nii sise- kui ka välistingimustes;
        vahetada patareid enne võistlust\\
        \hline
        Probleemid langevarjuga & Katsetada mitu korda; katsetada võimalikult suurelt kõrguselt; valmistada mitu langevarju mudelit\\
        \hline
        Probleemid korpusega & Katsetada korpust erinevatest materjalidest; valmistada korpuse varuosad\\
        \hline
        Ilma tingimused muutuvad & Valmistada ette mitu langevarju mudelit (erineva konstruktsioonidaga või pindalaga), katsetada neid ning valida kas universaalsema või erinevad erineva ilma tingimuste jaoks\\
        \hline
        Meeskonnaliikmed muutuvad & Dokumenteerida kõik oma poolt tehtud tööd, et saaks teda teisele liikmele üle anda; jälgida ajaplaani\\
        \hline
    \end{tabularx}
    \caption{Riskid ja vähendamise meetodid}
    \label{tab:4}
\end{table}

Peidetud riskid, mida ei osanud ette näha, kuid mis ilmnesid töö käigus:
\begin{itemize} 
    \item Meeskonnatöö raskused –-- oli keeruline kokku leppida mitmetes küsimustes: näiteks ühise tööaja leidmine, korpuse disaini ja omaduste osas kokkulepete saavutamine, aruande vormistamise detailid jms. Kõik need probleemid õnnestus lõpuks lahendada ja kokkulepped saavutada, kuid see nõudis lisatööd ja rohkem aega.
    \item Jaanuarist märtsini toimuvate olümpiaadide, võistluste ja konkursside periood –-- mitmed meeskonnaliikmed osalesid teistes projektides, mistõttu CanSati projektile jäi vähem aega või see tuli edasi lükata. Tulemuseks –-- kiirtöö (avarii töökorraldus) projekti lõpus.
    \item Ei arvestatud koolivaheaegu ega isiklikke sõite/sündmusi, mille ajal oli väga keeruline leida aega projekti tegevustega tegelemiseks.
    \item Ei arvestatud võimalikke \textit{force majeure}’i olukordi, mis samuti viivitasid meie projektitöö valmimist.
\end{itemize}

Järeldused:

\begin{itemize}
    \item Projekti ajaplaneerimisel tuleb arvestada võimalike ettenägematute olukordadega ning määrata tegevuste piirjooned varakult –-- enne projekti lõpptähtaegu.
    \item Ajaplaneerimisse tuleks lisada rohkem meeskonna füüsilisi kohtumisi, mille käigus saab ühiselt üle vaadata tehtud töö ja vormistada aruannet.
    \item Dokumentatsiooni on oluline koostada jooksvalt koos projektitegevustega --– nii on projekti lõpuks kõik sammud fikseeritud ning ei teki vajadust neid hiljem meelde tuletada ega uuesti taastada.

\end{itemize}
